\chapter{State of the Art}
\label{ch:estado-arte}

\section{Background}

This chapter frames malware analysis techniques for RAT detection and compares widely used
static, dynamic, and hybrid solutions with \rat.

\section{Malware Analysis: Overview}

Malware analysis is commonly divided into two complementary approaches
\cite{skoudis2004,cipriano2019,anderson2012,moser2007}. \textbf{Static analysis} examines the
file without executing it: PE structure, sections, imports/exports, strings, entropy, YARA
signatures, decompilation, and detection of evasion techniques (anti-debug, anti-VM).
\textbf{Dynamic analysis} runs the sample in an isolated sandbox and monitors behaviour---process
creation, registry changes, files, network activity, and persistence.

In the RAT context, static analysis identifies suspicious imports and C2 strings without
execution risk; dynamic analysis reveals connections, files, and persistence behaviour.
Neither approach is sufficient alone: packers and polymorphism defeat naive static signatures,
while sandbox evasion (anti-VM, sleep delays, environment checks) can suppress dynamic
indicators \cite{anderson2012,moser2007,bensaoud2024}.

The subsections below trace how these techniques evolved, where they fail in practice, and why
\rat\ favours a hybrid, explainable design.

\subsection{Evolution of Malware Analysis}

Malware analysis has evolved from manual disassembly and signature matching in the late 1990s
toward automated pipelines that combine sandbox execution, rule engines, and machine learning
\cite{skoudis2004,bayer2006,willems2007}. Early dynamic tools such as CWSandbox and TTAnalyze
demonstrated that behavioural traces could complement antivirus heuristics, but required
carefully controlled environments and human interpretation of noisy logs. The last decade brought
cloud multi-scanner services, open-source sandboxes (Cuckoo, CAPE), and deep-learning classifiers
on PE features \cite{raff2018,ucci2019,khan2024,maniriho2023}. Contemporary RAT families (AsyncRAT, QuasarRAT, Remcos)
further stress this landscape: they blend legitimate remote-administration APIs with obfuscation,
staged payloads, and commodity packers, making triage a \emph{multi-evidence} problem rather than
a single score.

\subsection{Static Versus Dynamic Detection: Challenges}

\textbf{Static analysis} is fast, repeatable, and safe---no execution on the analyst host---but
struggles with encrypted payloads, runtime decryption, and novel variants without known strings or
imports. \textbf{Dynamic analysis} observes ground-truth behaviour yet incurs higher latency, risks
sandbox escape if isolation fails, and may miss logic that activates only under specific user
interaction or after long delays \cite{cipriano2019,moser2007,alrawi2024}. Hybrid workflows therefore remain
the professional norm: static triage narrows the candidate set; dynamic confirmation validates
suspicious indicators \cite{bensaoud2024}. \rat\ adopts this hybrid model explicitly, prioritising explainable static
signals before optional sandbox execution.

\subsection{Explainability Versus Machine Learning}

Machine-learning detectors on raw bytes or PE feature vectors achieve strong benchmark metrics
\cite{raff2018,ucar2020,apruzzese2018,khan2024} but often function as opaque classifiers: an analyst
cannot easily justify a verdict to a student, auditor, or court. Recent work on explainable
behavioural detection \cite{galli2024} and interpretable Windows PE models \cite{khan2024,imran2024}
highlights this tension. In educational and laboratory settings, \textbf{interpretability} matters
as much as accuracy. Rule-based YARA matches, scored
heuristics (C2 strings, persistence, evasion), and decompiled pseudo-C provide auditable
evidence chains. Large language models offer a complementary layer---natural-language explanation
of code excerpts---without replacing the underlying indicators. \rat\ deliberately favours this
\emph{glass-box} design over end-to-end neural classification, accepting that novel malware may
score below threshold until rules or heuristics are updated.

\subsection{Limitations of Current Approaches}

No single commercial or open-source platform eliminates false positives and false negatives.
Multi-AV aggregation (VirusTotal) improves coverage but shares samples externally and hides
engine reasoning. Cloud sandboxes (Any.Run, Hybrid Analysis, Joe Sandbox) excel at rich telemetry
yet impose cost, latency, and confidentiality constraints for universities holding sensitive
student or research data. Self-hosted sandboxes (Cuckoo/CAPE) demand substantial operational
expertise. ML models require labelled datasets that are expensive to curate and drift as families
evolve \cite{vinod2019,mohaisen2016,imran2024}. These gaps motivate a \textbf{local}, \textbf{transparent},
and \textbf{modular} platform rather than another opaque cloud service.

\subsection{Sandbox Virtualisation: Hyper-V, Docker, and KVM}

Container platforms such as Docker are ill-suited to faithful Windows malware behaviour: many RATs
depend on Win32 APIs, GUI subsystems, registry persistence, and process injection patterns that
do not translate cleanly to Linux containers or minimal Windows images. \textbf{KVM/QEMU} offers
full hardware virtualisation and is common in Linux-centric research labs, but adds integration
overhead on Windows workstations where Hyper-V is already enabled for development. \textbf{Hyper-V}
provides type-1 hypervisor isolation on laboratory PCs, snapshot restore for repeatable runs, and
native PowerShell automation---trade-offs favouring a Windows-native academic deployment over
cross-platform portability. Docker and KVM remain valid alternatives where the host OS or
curriculum differs; they were evaluated and deferred (Section~\ref{ch:tecnologias}).

\subsection{Local Versus Cloud Analysis}

Submitting samples to public sandboxes accelerates triage but leaks file content to third parties
and prevents custom rule development inside the analysis environment. A \textbf{local sandbox}
keeps specimens on institutional infrastructure, supports offline teaching, and allows students to
inspect decompilation and scoring logic without account quotas. The cost is operational: the
analyst maintains VMs, snapshots, and updates. Configuration choices materially affect what a
sandbox observes \cite{alrawi2024}. \rat\ targets this local model consciously, accepting higher
per-run latency than commercial clouds in exchange for confidentiality, reproducibility, and
pedagogical depth.

\section{Trends in RAT and Malware Detection}

Beyond foundational static and dynamic methods, several product categories shape how analysts
detect RATs today. The following subsections review commercial sandboxes, multi-engine scanners,
open-source frameworks, and machine-learning approaches, highlighting trade-offs that informed the
design of \rat.

\subsection{Commercial and Cloud Sandboxes}

Platforms such as \emph{Any.Run} \cite{anyrun2024}, \emph{Hybrid Analysis} \cite{hybrid2019}, and
\emph{Joe Sandbox} \cite{joe2020} offer rich dynamic analysis \cite{bayer2006,willems2007,alrawi2024}, but
impose cost, external sample submission, and limited environment customisation---constraints that
\rat\ addresses with local, isolated, zero-cost execution.

\subsection{VirusTotal}

\emph{VirusTotal} \cite{virustotal2023} aggregates many antivirus engines for triage, but shares
submitted samples with partners and does not expose decompiled code or custom detection rules.

\subsection{Open-Source Frameworks}

Several open-source tools address parts of the malware-analysis pipeline, though none combine all
capabilities targeted by \rat\ in a single educational workflow. \textbf{Cuckoo Sandbox}
\cite{cuckoo2012} is a self-hostable, extensible Python sandbox with behavioural analysis on Windows
and Linux; it remains the open-source reference closest to \rat, but requires complex configuration
and does not integrate decompilation or explainable scoring. \textbf{CAPE Sandbox} \cite{cape2021}, a
Cuckoo fork, specialises in extracting configurations from malware families and C2 servers---highly
valuable, yet too narrow for generalist teaching. \textbf{YARA} \cite{yara2023} provides the rule
language used across the industry; \rat\ embeds YARA as one phase of the static pipeline alongside
scoring and decompilation. \textbf{Ghidra} and \textbf{ILSpy} \cite{ghidra2019,ilspy2024} are
essential decompilers typically used from the command line; \rat\ orchestrates both and exposes
results in the web interface.

\subsection{Machine Learning-Based Detection}

\emph{Deep learning} on PE features achieves high classification metrics
\cite{raff2018,ucar2020,ucci2019,apruzzese2018,maniriho2023} but yields opaque decisions. \rat\ prioritises
\textbf{explainability} via concrete indicators and an optional LLM assistant, aligning with recent
calls for transparent behavioural and static models \cite{galli2024,khan2024}. Public IOC
repositories such as MalwareBazaar \cite{abusech2021} and labelling tools such as AVclass
\cite{mohaisen2016} complement local triage \cite{vinod2019}.

\section{Positioning Relative to Existing Solutions}

\rat\ combines multi-phase static analysis with a native Hyper-V sandbox and integrated
decompilation, exposed through web and WPF interfaces. Table~\ref{tab:comparacao} summarises the
positioning.

\begin{table}[htbp]
    \centering
    \caption{Comparison of the \rat\ solution with existing alternatives}
    \label{tab:comparacao}
    \small
    \begin{tabularx}{\linewidth}{@{}>{\raggedright\arraybackslash}p{2.0cm}>{\raggedright\arraybackslash}p{4.0cm}Y@{}}
        \toprule
        \textbf{Reference} & \textbf{Technology} & \textbf{\rat\ solution} \\
        \midrule
        VirusTotal & Cloud multi-AV & Local; YARA + heuristics + Hyper-V sandbox \\
        Cuckoo/CAPE & Self-hostable sandbox & Native Windows sandbox (Hyper-V) \\
        Any.Run / Hybrid & Commercial sandbox & Free; educational focus; pseudo-C/IL \\
        YARA standalone & Signature rules & YARA + scoring + decompilation \\
        Ghidra/ILSpy CLI & Offline decompilation & Orchestrated decompilation + web UI \\
        ML classifiers & Opaque models & Explainable indicators + report + optional LLM assistant \\
        \bottomrule
    \end{tabularx}
\end{table}

\section{Conclusions}

Effective RAT detection benefits from multi-layer analysis that combines static indicators,
controlled execution, and human-interpretable reporting. The literature shows a clear tension
between scalable opaque ML and explainable rule-based triage; cloud sandboxes and local
hypervisor sandboxes serve different institutional needs. \rat\ positions itself for
\textbf{local academic use}: Hyper-V isolation, decompilation orchestration, scored heuristics,
and optional LLM-assisted code explanation, rather than commercial-scale multi-tenant analysis
or black-box classification alone.
